% =====================================================================
% Scholar 默认论文格式模板（FORMAT ONLY / 仅格式，不含写作风格）
% - 写稿流程：先写中文稿 -> 翻译英文 -> 套用本模板生成 .tex
% - 编译：python scripts/latex/compile_latex.py main.tex
% - 引擎：本模板用 ctex，自动走 xelatex（detect_engine 自动识别）
% - 说明：本模板只约束【格式】（版面/字体/章节/引用），
%         写作风格与内容规范见 standards/论文写作规范.md
% =====================================================================
\documentclass[11pt,a4paper]{article}

% ------------------------- 中文支持（xelatex） -------------------------
\usepackage{ctex}          % 中文稿编译预览；最终英文稿可保留也可移除
\usepackage{fontspec}

% ------------------------- 页面与字体 -------------------------
\usepackage[a4paper, margin=1in]{geometry}
\usepackage{microtype}
\usepackage{setspace}
\onehalfspacing

% ------------------------- 数学 -------------------------
\usepackage{amsmath}
\usepackage{amssymb}
\usepackage{mathtools}

% ------------------------- 图表 -------------------------
\usepackage{graphicx}
\usepackage[font=small,labelfont=bf]{caption}
\usepackage{subcaption}
\usepackage{booktabs}
\usepackage{array}
\usepackage{multirow}

% ------------------------- 列表 -------------------------
\usepackage{enumitem}
\setlist{nosep}

% ------------------------- 颜色与超链接 -------------------------
\usepackage{xcolor}
\definecolor{linkblue}{RGB}{0,51,153}
\usepackage[colorlinks=true,linkcolor=linkblue,citecolor=linkblue,urlcolor=linkblue]{hyperref}

% ------------------------- 算法 -------------------------
\usepackage{algorithm}
\usepackage{algpseudocode}

% ------------------------- 参考文献 -------------------------
\usepackage{natbib}
\bibliographystyle{plainnat}

% =====================================================================
\begin{document}

\title{Paper Title}
\author{Author One$^{1}$, Author Two$^{2}$ \\ \small $^{1}$Affiliation 1, $^{2}$Affiliation 2}
\date{\today}
\maketitle

% ------------------------- 摘要 -------------------------
\begin{abstract}
% 中文稿阶段可写中文摘要；英文稿阶段替换为英文摘要。
Abstract text: background, problem, method, key results, conclusion.
\end{abstract}

\noindent\textbf{Keywords:} keyword1; keyword2; keyword3

% ------------------------- 正文 -------------------------
\section{Introduction}
% 背景、动机、问题、贡献列表、论文组织。

\section{Related Work}
% 按主题组织，避免逐篇枚举；每主题末尾给批判性小结。

\section{Method}
% 问题形式化、算法设计、数学与理论；伪代码用 algorithm 环境。

\section{Experiments}
% 数据集/设置、基线、实现细节；对照解决方案文档的评价指标。

\section{Results}
% 主结果、消融、鲁棒性；图表引用 \ref{}，指标表格用 booktabs。

\section{Conclusion}
% 结论、局限、未来工作。

% ------------------------- 参考文献 -------------------------
\bibliography{references}

\end{document}
